%Metamath.tex

\section{Metamath Preliminaries}\label{sec:metamathprelim}

\bemph{Metamath}\footnote{http://us.metamath.org/index.html} will be used to prove as theorems in set theory, what are assumed by the BLESS proof tool as axioms and inference rules.  Metamath is a
\begin{itemize}
  \item language to express mathematics, thus ``meta-''
  \item tool suite; especially Metamath Proof Explorer (checker)
  \item library of over 10,000 proofs, individually named and consecutively numbered, starting from set theory axioms 
\end{itemize}

%Metamath is hard to understand.  XML files include ``metadata", data-about-data, to flexibly-express structured information.  The XML metadata defines the structure of the information.  

The ``meta-'' in Metamath means its subject is math.  Metamath does math about math.  

Metamath allows you to define your own mathematical system.  Once defined you can prove interesting things about your mathematical system.

Metamath language allows expression of (other) mathematical systems, and prove theorems about those mathematical systems, which will in turn be used to prove theorems about the math you want to use.  Here Metamath will used to prove soundness about the proof rules used by the BLESS proof engine.  For BLESS, Metamath is used to prove that each proof rule is \bemph{sound}, either tautology or inference rule, always producing true facts from true facts.

\section{Metamath Symbols}\label{sec:metamathsymbols}

Metamath symbols (Table \ref{tab:sym}) follow customary usage.  Symbols introduced especially for BLESS (Table \ref{tab:blesssym}) are explained when their use is defined.  Variable symbols (Table \ref{tab:var}) have what are effectively types.

Most of the steps in a Metamath proof assure the syntax is correct.  As part of this, variable symbols can be used in ways such that the syntax enforces types.  The proofs pretty-printed in \LaTeX omit steps that assure syntactical correctness.

%For example, the proof for bl.that (\ref{}) is represented in the bless.mm database as
%\begin{lstlisting}
%${
%bl.that.1 $e |- ph $.
%$( theorems are true $)
%bl.that $p |- ( ph <-> true )
%	$=
%	$( wffs for bitr4i  ph= ph  ps= ( ph -> ph )  ch= true  $)
%	  wph wph wph wi wtrue
%	$( wffs for a1bi ph= ph  ps= ph $)
%	  wph wph
%	bl.that.1
%	a1bi  $( |- ( ph <-> ( ph -> ph ) )  $)
%	wph df-true $( |- ( true <-> ( ph -> ph ) ) $)
%	$( by bitr4i |- ( ph <-> ps )  and  |- ( ch <-> ps )  =>  |- ( ph <-> ch )
%		with ph= ph  ps= ( ph -> ph )  ch= true  $)
%	bitr4i
%	$( |- ( ph <-> true ) $)
%	$.
%$}
%$( end of bl.that $)
%\end{lstlisting}
%
%Symbols \lstba{$(} and \lstba{$)} delimit comments;  \lstba{$\{} and \lstba{$\}} delimit scope.

%\begin{description}
%\item[$$] 
%\end{description}

%% This LaTeX file was created by Metamath on 4-Nov-2016 3:06 PM.
\begin{lemma}\otherTitle{wsbw} Substitution of a wff by a wff in a wff
\begin{align}
{\rm\color{wffcolor} wff~} \overline{[} {\color{formulacolor}\psi} /
    {\color{formulacolor}\chi} \overline{]}
    {\color{formulacolor}\varphi}\label{wsbw}
\end{align}
\end{lemma}

\begin{lemma}\otherTitle{ccsl1} declare 'csl1' to be a csl
\begin{align}
{\rm\color{wffcolor} csl}~ {\rm\color{variablecolor}csl_1}\label{ccsl1}
\end{align}
\end{lemma}

\begin{lemma}\otherTitle{ccsl2} declare 'csl2' to be a csl
\begin{align}
{\rm\color{wffcolor} csl}~ {\rm\color{variablecolor}csl_2}\label{ccsl2}
\end{align}
\end{lemma}

\begin{lemma}\otherTitle{as} declare 'S' to be an action
\begin{align}
{\rm\color{wffcolor}action}~ {\color{actioncolor}S}\label{as}
\end{align}
\end{lemma}

\begin{lemma}\otherTitle{as1} declare 'S1' to be an action
\begin{align}
{\rm\color{wffcolor}action}~ {\color{actioncolor}S_1}\label{as1}
\end{align}
\end{lemma}

\begin{lemma}\otherTitle{as2} declare 'S2' to be an action
\begin{align}
{\rm\color{wffcolor}action}~ {\color{actioncolor}S_2}\label{as2}
\end{align}
\end{lemma}

\begin{lemma}\otherTitle{as3} declare 'S3' to be an action
\begin{align}
{\rm\color{wffcolor}action}~ {\color{actioncolor}S_3}\label{as3}
\end{align}
\end{lemma}

\begin{lemma}\otherTitle{asn} declare 'Sn' to be an action
\begin{align}
{\rm\color{wffcolor}action}~ {\color{actioncolor}S_n}\label{asn}
\end{align}
\end{lemma}

\begin{lemma}\otherTitle{assP} declare 'P-' to be an assertion
\begin{align}
{\rm\color{wffcolor}assert}~ {\color{assertioncolor}P}\label{assP}
\end{align}
\end{lemma}

\begin{lemma}\otherTitle{assP1} TO DO: PUT DESCRIPTION HERE
\begin{align}
{\rm\color{wffcolor}assert}~ {\color{assertioncolor}P_1}\label{assP1}
\end{align}
\end{lemma}

\begin{lemma}\otherTitle{assP2} TO DO: PUT DESCRIPTION HERE
\begin{align}
{\rm\color{wffcolor}assert}~ {\color{assertioncolor}P_2}\label{assP2}
\end{align}
\end{lemma}

\begin{lemma}\otherTitle{assP3} TO DO: PUT DESCRIPTION HERE
\begin{align}
{\rm\color{wffcolor}assert}~ {\color{assertioncolor}P_3}\label{assP3}
\end{align}
\end{lemma}

\begin{lemma}\otherTitle{assPn} TO DO: PUT DESCRIPTION HERE
\begin{align}
{\rm\color{wffcolor}assert}~ {\color{assertioncolor}P_n}\label{assPn}
\end{align}
\end{lemma}

\begin{lemma}\otherTitle{assQ} declare 'Q-' to be an assertion
\begin{align}
{\rm\color{wffcolor}assert}~ {\color{assertioncolor}Q}\label{assQ}
\end{align}
\end{lemma}

\begin{lemma}\otherTitle{assQ1} TO DO: PUT DESCRIPTION HERE
\begin{align}
{\rm\color{wffcolor}assert}~ {\color{assertioncolor}Q_1}\label{assQ1}
\end{align}
\end{lemma}

\begin{lemma}\otherTitle{assQ2} TO DO: PUT DESCRIPTION HERE
\begin{align}
{\rm\color{wffcolor}assert}~ {\color{assertioncolor}Q_2}\label{assQ2}
\end{align}
\end{lemma}

\begin{lemma}\otherTitle{assQ3} TO DO: PUT DESCRIPTION HERE
\begin{align}
{\rm\color{wffcolor}assert}~ {\color{assertioncolor}Q_3}\label{assQ3}
\end{align}
\end{lemma}

\begin{lemma}\otherTitle{assQn} TO DO: PUT DESCRIPTION HERE
\begin{align}
{\rm\color{wffcolor}assert}~ {\color{assertioncolor}Q_n}\label{assQn}
\end{align}
\end{lemma}

\begin{lemma}\otherTitle{aa1} TO DO: PUT DESCRIPTION HERE
\begin{align}
{\rm\color{wffcolor}aa}~ A_1\label{aa1}
\end{align}
\end{lemma}

\begin{lemma}\otherTitle{aa2} TO DO: PUT DESCRIPTION HERE
\begin{align}
{\rm\color{wffcolor}aa}~ A_2\label{aa2}
\end{align}
\end{lemma}

\begin{lemma}\otherTitle{aa3} TO DO: PUT DESCRIPTION HERE
\begin{align}
{\rm\color{wffcolor}aa}~ A_3\label{aa3}
\end{align}
\end{lemma}

\begin{lemma}\otherTitle{aan} TO DO: PUT DESCRIPTION HERE
\begin{align}
{\rm\color{wffcolor}aa}~ A_n\label{aan}
\end{align}
\end{lemma}

\begin{lemma}\otherTitle{wl1} declare 'l-' to be a wl
\begin{align}
{\rm\color{wffcolor} wl}~ {\rm\color{variablecolor}l_1}\label{wl1}
\end{align}
\end{lemma}

\begin{lemma}\otherTitle{wl2} declare 'l2' to be a wl
\begin{align}
{\rm\color{wffcolor} wl}~ {\rm\color{variablecolor}l_2}\label{wl2}
\end{align}
\end{lemma}

\begin{lemma}\otherTitle{wl3} declare 'l3' to be a wl
\begin{align}
{\rm\color{wffcolor} wl}~ {\rm\color{variablecolor}l_3}\label{wl3}
\end{align}
\end{lemma}

\begin{lemma}\otherTitle{wl4} declare 'l4' to be a wl
\begin{align}
{\rm\color{wffcolor} wl}~ {\rm\color{variablecolor}l_4}\label{wl4}
\end{align}
\end{lemma}

\begin{lemma}\otherTitle{wl5} declare 'l5' to be a wl
\begin{align}
{\rm\color{wffcolor} wl}~ {\rm\color{variablecolor}l_5}\label{wl5}
\end{align}
\end{lemma}

\begin{lemma}\otherTitle{csl-a} Define Comma-Separated List: one element A
\begin{align}
{\rm\color{wffcolor} csl}~ {\color{classcolor}A}\label{csl-a}
\end{align}
\end{lemma}

\begin{lemma}\otherTitle{csl-acl} problem generating LaTeX for csl-acl
\begin{align}
{\rm\color{wffcolor} csl}~ {\rm\color{variablecolor}csl_1} ,
    {\color{classcolor}A}\label{csl-acl}
\end{align}
\end{lemma}

\begin{lemma}\otherTitle{cl-ntupl} Extend class notation to include  N-tuple
\begin{align}
{\rm class} \langle {\rm\color{variablecolor}csl_1} \rangle\label{cl-ntupl}
\end{align}
\end{lemma}

\begin{lemma}\otherTitle{wfl0} wff-list zero wffs
\begin{align}
{\rm\color{wffcolor} wl}~\label{wfl0}
\end{align}
\end{lemma}

\begin{lemma}\otherTitle{wfl1} wff-list one wff
\begin{align}
{\rm\color{wffcolor} wl}~ {\color{formulacolor}\varphi}\label{wfl1}
\end{align}
\end{lemma}

\begin{lemma}\otherTitle{wfl2} two-element wff-list
\begin{align}
{\rm\color{wffcolor} wl}~ {\rm\color{variablecolor}l_1}
    {\rm\color{variablecolor}l_2}\label{wfl2}
\end{align}
\end{lemma}

\begin{lemma}\otherTitle{wland} Multi-Term And:  define conjunction of a list
of wff
\begin{align}
{\rm\color{wffcolor} wff~} \bigwedge( {\rm\color{variablecolor}l_1}
    )\label{wland}
\end{align}
\end{lemma}

\begin{lemma}\otherTitle{wlo} Multi-Term Or:  define disjunction of a list of
wff
\begin{align}
{\rm\color{wffcolor} wff~} \bigvee( {\rm\color{variablecolor}l_1}
    )\label{wlo}
\end{align}
\end{lemma}

\begin{lemma}\otherTitle{assrt0} an assertion is predicate between angle
brackets
\begin{align}
{\rm\color{wffcolor}assert}~ \llx {\color{formulacolor}\varphi}
    \xgg\label{assrt0}
\end{align}
\end{lemma}

\begin{lemma}\otherTitle{aaarrow} make \{\{. ph \}\}. -\} \{\{. ps \}\}.
asserted action
\begin{align}
{\rm\color{wffcolor}aa}~ {\color{assertioncolor}P} \rightarrow
    {\color{assertioncolor}Q}\label{aaarrow}
\end{align}
\end{lemma}

\begin{lemma}\otherTitle{aapsq} make \{\{. ph \}\}. S$_\{1\}$ \{\{. ps \}\}.
asserted action
\begin{align}
{\rm\color{wffcolor}aa}~ {\color{assertioncolor}P} {\color{actioncolor}S_1}
    {\color{assertioncolor}Q}\label{aapsq}
\end{align}
\end{lemma}

\begin{lemma}\otherTitle{aaps} make \{\{. ph \}\}. S1 asserted action
\begin{align}
{\rm\color{wffcolor}aa}~ {\color{assertioncolor}P}
    {\color{actioncolor}S_1}\label{aaps}
\end{align}
\end{lemma}

\begin{lemma}\otherTitle{aasq} make S$_\{1\}$ \{\{. ps \}\}. asserted action
\begin{align}
{\rm\color{wffcolor}aa}~ {\color{actioncolor}S_1}
    {\color{assertioncolor}Q}\label{aasq}
\end{align}
\end{lemma}

\begin{lemma}\otherTitle{aas} make  S$_\{1\}$ asserted action
\begin{align}
{\rm\color{wffcolor}aa}~ {\color{actioncolor}S_1}\label{aas}
\end{align}
\end{lemma}

\begin{lemma}\otherTitle{wtrue} true is a well-formed formula
\begin{align}
{\rm\color{wffcolor} wff~} \top\label{wtrue}
\end{align}
\end{lemma}

\begin{lemma}\otherTitle{wfalse} false is a well-formed formula
\begin{align}
{\rm\color{wffcolor} wff~} \bot\label{wfalse}
\end{align}
\end{lemma}



\pagebreak
 \section{Axioms}\label{sec:axioms}

\pn The axioms of \texttt{set.mm} are listed below.  The letters in parentheses by the right-hand margin is the name of the axiom when used in proofs.

\pn The first `axiom' is really the inference rule, Modus Ponens.  It's called `axiom' by Metamath because it's assumed.  Other inference rules will be proved using Modus Ponens as theorems.  The $\vdash$ symbol means ``a proof exists for", or more colloquially ``is true"..  The $\&$ and $\Rightarrow$ symbols are part of the simple metalanguage used by Metamath, and are not used in theorems or proofs.  $\&$ separates hypotheses of an inference rule\footnote{BLESS will introduce a different ampersand that means concurrent composition.}; $\Rightarrow$ separates the hypotheses from the conclusion.

\pn The Rule of \bemph{Modus Ponens} states:  if a proof exists for $\varphi$, and a proof exists that $\varphi$ implies $\psi$ ($\rightarrow$ means implies), then we can conclude that $\psi$ has a proof.

\begin{axiom}\label{axi:ax-mp} Rule of Modus Ponens.  
\begin{align}
\vdash \varphi\quad\&\quad \vdash ( \varphi \rightarrow \psi
    )\quad\Rightarrow\quad \vdash \psi\label{eq:ax-mp}\tag{ax-mp}
\end{align}
\end{axiom}

\subsection{Axioms of Propositional Logic}\label{sec:proplogic}

\pn Propositions are statements which are either true or false, if not meaningless.  "This statement is a lie." is an example of a meaningless statement that is neither true nor false. $\lnot$ means not.

\begin{axiom}\label{axi:ax-1} Axiom of Simplification 
\begin{align}
\vdash ( \varphi \rightarrow ( \psi \rightarrow \varphi )
    )\label{eq:ax-1}\tag{ax-1}
\end{align}
\end{axiom}
\begin{axiom}\label{axi:ax-2} Axiom of Distribution 
\begin{align}
\vdash ( ( \varphi \rightarrow ( \psi \rightarrow \chi ) ) \rightarrow ( (
    \varphi \rightarrow \psi ) \rightarrow ( \varphi \rightarrow \chi ) )
    )\label{eq:ax-2}\tag{ax-2}
\end{align}
\end{axiom}
\begin{axiom}\label{axi:ax-3} Axiom of Contraposition   
\begin{align}
\vdash ( ( \lnot \varphi \rightarrow \lnot \psi ) \rightarrow ( \psi
    \rightarrow \varphi ) )\label{eq:ax-3}\tag{ax-3}
\end{align}
\end{axiom}

\subsection{Axioms of Predicate Calculus with Equality}

\pn Predicate calculus extends predicate logic with variables, quantifiers, and equality.  $\forall$ means for-all;
$\exists$ means there-exists. 

\begin{axiom}\label{axi:ax-gen} Rule of Generalization. 
\begin{align}
\vdash \varphi\quad\Rightarrow\quad \vdash \forall x
    \varphi\label{eq:ax-gen}\tag{ax-gen}
\end{align}
\end{axiom}
\begin{axiom}\label{axi:ax-4} Axiom of Quantified Implication. 
\begin{align}
\vdash ( \forall x ( \varphi \rightarrow \psi ) \rightarrow ( \forall x \varphi
    \rightarrow \forall x \psi ) )\label{eq:ax-4}\tag{ax-4}
\end{align}
\end{axiom}
\begin{axiom}\label{axi:ax-5} Axiom of Distinctness.  \begin{align}
\vdash ( \varphi \rightarrow \forall x \varphi )\label{eq:ax-5}\tag{ax-5}
\end{align}
\end{axiom}
\begin{axiom}\label{axi:ax-6} Axiom of Existence. 
\begin{align}
\vdash \lnot \forall x \lnot x = y\label{eq:ax-6}\tag{ax-6}
\end{align}
\end{axiom}
\begin{axiom}\label{axi:ax-7} Axiom of Equality.  
\begin{align}
\vdash ( x = y \rightarrow ( x = z \rightarrow y = z )
    )\label{eq:ax-7}\tag{ax-7}
\end{align}
\end{axiom}
\begin{axiom}\label{axi:ax-8} Axiom of Left Equality for Binary Predicate. 
\begin{align}
\vdash ( x = y \rightarrow ( x \in z \rightarrow y \in z )
    )\label{eq:ax-8}\tag{ax-8}
\end{align}
\end{axiom}
\begin{axiom}\label{axi:ax-9} Axiom of Right Equality for Binary Predicate. 
\begin{align}
\vdash ( x = y \rightarrow ( z \in x \rightarrow z \in y )
    )\label{eq:ax-9}\tag{ax-9}
\end{align}
\end{axiom}
\begin{axiom}\label{axi:ax-10} Axiom of Quantified Negation.  
\begin{align}
\vdash ( \lnot \forall x \varphi \rightarrow \forall x \lnot \forall x \varphi
    )\label{eq:ax-10}\tag{ax-10}
\end{align}
\end{axiom}
\begin{axiom}\label{axi:ax-11} Axiom of Quantifier Commutation.  
\begin{align}
\vdash ( \forall x \forall y \varphi \rightarrow \forall y \forall x \varphi
    )\label{eq:ax-11}\tag{ax-11}
\end{align}
\end{axiom}
\begin{axiom}\label{axi:ax-12} Axiom of Substitution.  
\begin{align}
\vdash ( x = y \rightarrow ( \forall y \varphi \rightarrow \forall x ( x = y
    \rightarrow \varphi ) ) )\label{eq:ax-12}\tag{ax-12}
\end{align}
\end{axiom}
\begin{axiom}\label{axi:ax-13} Axiom of Quantified Equality. 
\begin{align}
\vdash ( \lnot x = y \rightarrow ( y = z \rightarrow \forall x \,y = z )
    )\label{eq:ax-13}\tag{ax-13}
\end{align}
\end{axiom}

\subsection{Axioms of Zermelo-Frankel with Choice (ZFC) set theory}\label{sec.zfc}

\pn Predicate calculus introduced variables and the equality relation between them, but wasn't specific
about what they were.  With the addition of axioms of set theory, variables are sets.

\begin{axiom}\label{axi:ax-ext} Axiom of Extensionality.  
\begin{align}
\vdash ( \forall z ( z \in x \leftrightarrow z \in y ) \rightarrow x = y
    )\label{eq:ax-ext}\tag{ax-ext}
\end{align}
\end{axiom}
\begin{axiom}\label{axi:ax-rep} Axiom of Replacement.  
\begin{align}
\vdash ( \forall w \exists y \forall z ( \forall y \varphi \rightarrow z = y )
    \rightarrow \exists y \forall z ( z \in y \leftrightarrow \exists w ( w \in
    x \wedge \forall y \varphi ) ) )\label{eq:ax-rep}\tag{ax-rep}
\end{align}
\end{axiom}
\begin{axiom}\label{axi:ax-pow} Axiom of Power Sets. 
\begin{align}
\vdash \exists y \forall z ( \forall w ( w \in z \rightarrow w \in x )
    \rightarrow z \in y )\label{eq:ax-pow}\tag{ax-pow}
\end{align}
\end{axiom}
\begin{axiom}\label{axi:ax-un} Axiom of Union. 
\begin{align}
\vdash \exists y \forall z ( \exists w ( z \in w \wedge w \in x ) \rightarrow z
    \in y )\label{eq:ax-un}\tag{ax-un}
\end{align}
\end{axiom}
\begin{axiom}\label{axi:ax-reg} Axiom of Regularity. 
\begin{align}
\vdash ( \exists y \,y \in x \rightarrow \exists y ( y \in x \wedge \forall z (
    z \in y \rightarrow \lnot z \in x ) ) )\label{eq:ax-reg}\tag{ax-reg}
\end{align}
\end{axiom}
\begin{axiom}\label{axi:ax-inf} Axiom of Infinity. 
\begin{align}
\vdash \exists y ( x \in y \wedge \forall z ( z \in y \rightarrow \exists w ( z
    \in w \wedge w \in y ) ) )\label{eq:ax-inf}\tag{ax-inf}
\end{align}
\end{axiom}
\begin{axiom}\label{axi:ax-ac} Axiom of Choice.  
\begin{align}
\vdash \exists y \forall z \forall w ( ( z \in w \wedge w \in x ) \rightarrow 
    \exists v \forall u ( \exists t ( ( u \in w \wedge w \in t ) \wedge ( u \in
    t \wedge t \in y ) ) \leftrightarrow u = v ) )\label{eq:ax-ac}\tag{ax-ac}
\end{align}
\end{axiom}


% \section{Metamath Definitions Theorems and Axioms}\label{sec:tanda}
%
%\begin{axiom}\label{axi:ax-mp} Rule of Modus Ponens.  
%\begin{align}
%\vdash \varphi\quad\&\quad \vdash ( \varphi \rightarrow \mw{\psi}
%    )\quad\Rightarrow\quad \vdash \mw{\psi}\label{eq:ax-mp}\tag{ax-mp}
%\end{align}
%\end{axiom}
%\begin{axiom}\label{axi:ax-1} Axiom of Simplification 
%\begin{align}
%\vdash ( \varphi \rightarrow ( \mw{\psi} \rightarrow \varphi )
%    )\label{eq:ax-1}\tag{ax-1}
%\end{align}
%\end{axiom}
%\begin{axiom}\label{axi:ax-2} Axiom of Distribution 
%\begin{align}
%\vdash ( ( \varphi \rightarrow ( \mw{\psi} \rightarrow \chi ) ) \rightarrow ( (
%    \varphi \rightarrow \mw{\psi} ) \rightarrow ( \varphi \rightarrow \chi ) )
%    )\label{eq:ax-2}\tag{ax-2}
%\end{align}
%\end{axiom}
%\begin{axiom}\label{axi:ax-3} Axiom of Contraposition   
%\begin{align}
%\vdash ( ( \lnot \varphi \rightarrow \lnot \mw{\psi} ) \rightarrow ( \mw{\psi}
%    \rightarrow \varphi ) )\label{eq:ax-3}\tag{ax-3}
%\end{align}
%\end{axiom}
%\begin{axiom}\label{axi:ax-gen} Rule of Generalization. 
%\begin{align}
%\vdash \varphi\quad\Rightarrow\quad \vdash \forall x
%    \varphi\label{eq:ax-gen}\tag{ax-gen}
%\end{align}
%\end{axiom}
%\begin{axiom}\label{axi:ax-4} Axiom of Quantified Implication. 
%\begin{align}
%\vdash ( \forall x ( \varphi \rightarrow \mw{\psi} ) \rightarrow ( \forall x \varphi
%    \rightarrow \forall x \mw{\psi} ) )\label{eq:ax-4}\tag{ax-4}
%\end{align}
%\end{axiom}
%\begin{axiom}\label{axi:ax-5} Axiom of Distinctness.  \begin{align}
%\vdash ( \varphi \rightarrow \forall x \varphi )\label{eq:ax-5}\tag{ax-5}
%\end{align}
%\end{axiom}
%\begin{axiom}\label{axi:ax-6} Axiom of Existence. 
%\begin{align}
%\vdash \lnot \forall x \lnot x = y\label{eq:ax-6}\tag{ax-6}
%\end{align}
%\end{axiom}
%\begin{axiom}\label{axi:ax-7} Axiom of Equality.  
%\begin{align}
%\vdash ( x = y \rightarrow ( x = z \rightarrow y = z )
%    )\label{eq:ax-7}\tag{ax-7}
%\end{align}
%\end{axiom}
%\begin{axiom}\label{axi:ax-8} Axiom of Left Equality for Binary Predicate. 
%\begin{align}
%\vdash ( x = y \rightarrow ( x \in z \rightarrow y \in z )
%    )\label{eq:ax-8}\tag{ax-8}
%\end{align}
%\end{axiom}
%\begin{axiom}\label{axi:ax-9} Axiom of Right Equality for Binary Predicate. 
%\begin{align}
%\vdash ( x = y \rightarrow ( z \in x \rightarrow z \in y )
%    )\label{eq:ax-9}\tag{ax-9}
%\end{align}
%\end{axiom}
%\begin{axiom}\label{axi:ax-10} Axiom of Quantified Negation.  
%\begin{align}
%\vdash ( \lnot \forall x \varphi \rightarrow \forall x \lnot \forall x \varphi
%    )\label{eq:ax-10}\tag{ax-10}
%\end{align}
%\end{axiom}
%\begin{axiom}\label{axi:ax-11} Axiom of Quantifier Commutation.  
%\begin{align}
%\vdash ( \forall x \forall y \varphi \rightarrow \forall y \forall x \varphi
%    )\label{eq:ax-11}\tag{ax-11}
%\end{align}
%\end{axiom}
%\begin{axiom}\label{axi:ax-12} Axiom of Substitution.  
%\begin{align}
%\vdash ( x = y \rightarrow ( \forall y \varphi \rightarrow \forall x ( x = y
%    \rightarrow \varphi ) ) )\label{eq:ax-12}\tag{ax-12}
%\end{align}
%\end{axiom}
%\begin{axiom}\label{axi:ax-13} Axiom of Quantified Equality. 
%\begin{align}
%\vdash ( \lnot x = y \rightarrow ( y = z \rightarrow \forall x \,y = z )
%    )\label{eq:ax-13}\tag{ax-13}
%\end{align}
%\end{axiom}
%\begin{axiom}\label{axi:ax-ext} Axiom of Extensionality.  
%\begin{align}
%\vdash ( \forall z ( z \in x \leftrightarrow z \in y ) \rightarrow x = y
%    )\label{eq:ax-ext}\tag{ax-ext}
%\end{align}
%\end{axiom}
%\begin{axiom}\label{axi:ax-rep} Axiom of Replacement.  
%\begin{align}
%\vdash ( \forall w \exists y \forall z ( \forall y \varphi \rightarrow z = y )
%    \rightarrow \exists y \forall z ( z \in y \leftrightarrow \exists w ( w \in
%    x \wedge \forall y \varphi ) ) )\label{eq:ax-rep}\tag{ax-rep}
%\end{align}
%\end{axiom}
%\begin{axiom}\label{axi:ax-pow} Axiom of Power Sets. 
%\begin{align}
%\vdash \exists y \forall z ( \forall w ( w \in z \rightarrow w \in x )
%    \rightarrow z \in y )\label{eq:ax-pow}\tag{ax-pow}
%\end{align}
%\end{axiom}
%\begin{axiom}\label{axi:ax-un} Axiom of Union. 
%\begin{align}
%\vdash \exists y \forall z ( \exists w ( z \in w \wedge w \in x ) \rightarrow z
%    \in y )\label{eq:ax-un}\tag{ax-un}
%\end{align}
%\end{axiom}
%\begin{axiom}\label{axi:ax-reg} Axiom of Regularity. 
%\begin{align}
%\vdash ( \exists y \,y \in x \rightarrow \exists y ( y \in x \wedge \forall z (
%    z \in y \rightarrow \lnot z \in x ) ) )\label{eq:ax-reg}\tag{ax-reg}
%\end{align}
%\end{axiom}
%\begin{axiom}\label{axi:ax-inf} Axiom of Infinity. 
%\begin{align}
%\vdash \exists y ( x \in y \wedge \forall z ( z \in y \rightarrow \exists w ( z
%    \in w \wedge w \in y ) ) )\label{eq:ax-inf}\tag{ax-inf}
%\end{align}
%\end{axiom}
%\begin{axiom}\label{axi:ax-ac} Axiom of Choice.  
%\begin{align}
%\vdash \exists y \forall z \forall w ( ( z \in w \wedge w \in x ) \rightarrow 
%    \exists v \forall u ( \exists t ( ( u \in w \wedge w \in t ) \wedge ( u \in
%    t \wedge t \in y ) ) \leftrightarrow u = v ) )\label{eq:ax-ac}\tag{ax-ac}
%\end{align}
%\end{axiom}

\subsection{Class and Well-formed Formula Definitions}\label{sec:clwff}

\begin{definition}\otherTitle{cv} Set variable is also class variable
\begin{align}
{\rm class} x\label{eq:cv}\tag{cv}
\end{align}
\end{definition}
\begin{definition}\otherTitle{wcel} Extend wff definition to include the
membership connective between
       classes.
\begin{align}
{\rm wff} A \in B\label{eq:wcel}\tag{wcel}
\end{align}
\end{definition}

\begin{definition}\otherTitle{wa} Extend wff definition to include conjunction
('and').
\begin{align}
{\rm wff} ( \varphi \wedge \mw{\psi} )\label{eq:wa}\tag{wa}
\end{align}
\end{definition}
\begin{definition}\otherTitle{wi} If $\varphi$ and $\psi$ are wff's, so is $(
\varphi  \rightarrow  \mw{\psi} )$ or `` $\varphi$ implies
     $\psi$." 
\begin{align}
{\rm wff} ( \varphi \rightarrow \mw{\psi} )\label{eq:wi}\tag{wi}
\end{align}
\end{definition}

\subsection{Symbol Definitions}\label{sec:symdef}

\begin{definition}\otherTitle{df-an} Define conjunction (logical 'and'). 
\begin{align}
\vdash ( ( \varphi \wedge \mw{\psi} ) \leftrightarrow \lnot ( \varphi \rightarrow
    \lnot \mw{\psi} ) )\label{eq:df-an}\tag{df-an}
\end{align}
\end{definition}

\begin{definition}\otherTitle{df-or} Define disjunction (logical 'or'). 
\begin{align}
\vdash ( ( \varphi \vee \mw{\psi} ) \leftrightarrow ( \lnot \varphi \rightarrow \mw{\psi}
    ) )\label{eq:df-or}\tag{df-or}
\end{align}
\end{definition}

\begin{definition}\otherTitle{df-ltxr} Define 'less than' on the set of
extended reals. 
\begin{align}
\vdash < = ( \{ \langle x , y \rangle | ( x \in \mathbb{R} \wedge y \in
    \mathbb{R} \wedge x <_\mathbb{R} y ) \} \cup ( ( ( \mathbb{R} \cup \{
    -\infty \} ) \times \{ +\infty \} ) \cup ( \{ -\infty \} \times \mathbb{R}
    ) ) )\label{eq:df-ltxr}\tag{df-ltxr}
\end{align}
\end{definition}

\begin{definition}\otherTitle{df-tru} Definition of the truth value ``true", or
``verum", denoted by $\top$.
\begin{align}
\vdash ( \top \leftrightarrow ( \forall x \,x = x \rightarrow \forall x \,x = x
    ) )\label{eq:df-tru}\tag{df-tru}
\end{align}

\end{definition}
\begin{definition}\otherTitle{df-fal} Definition of the truth value ``false",
or ``falsum", denoted by $\bot$.
\begin{align}
\vdash ( \bot \leftrightarrow \lnot \top )\label{eq:df-fal}\tag{df-fal}
\end{align}
\end{definition}



%\subsection{Ordinal Numbers}\label{sec:ordinal}
%
%\pn Ordinal numbers create natural numbers out of pure set theory.  Just the existence of The Empty Set, and the operation of creating a new set by listing its elements.
%
%\pn I'm rather skeptical about the existence of The Empty Set, but accept the hypothesis for now.
%
%% This LaTeX file was created by Metamath on 9-Oct-2021 9:09 AM.
% hand edited by brl

The Empty Set is an ordinal class {Ord}, and it's a member of the set of ordinals (On)

\begin{lemma}\blTitle{ord0} The empty set is an ordinal class. 
\begin{align}
\mm{\vdash} {\rm Ord}~ \varnothing\label{eq:ord0}\tag{ord0}
\end{align}
\end{lemma}

\begin{lemma}\blTitle{0elon} The empty set is an ordinal number. 
\begin{align}
\mm{\vdash} \varnothing \in {\rm On}\label{eq:0elon}\tag{0elon}
\end{align}
\end{lemma}


\pn Then use the successor operation to make as many ordinal numbers as you want.

\begin{metadefinition}\blTitle{df-suc} Define the successor of a class.  When
applied to an ordinal number, the
     successor means the same thing as ``plus 1" 
\begin{align}
\mm{\vdash} {\rm suc}~ \mc{A} = ( \mc{A} \cup \{ \mc{A} \}
    )\label{eq:df-suc}\tag{df-suc}
\end{align}
\end{metadefinition}


\pn If $\varnothing$ represents 0, then its successor should represent 1.

\begin{lemma}\blTitle{suc0} The successor of the empty set.  (Contributed
by NM, 1-Feb-2005.)
\begin{align}
\mm{\vdash} {\rm suc}~ \varnothing = \{ \varnothing \}\label{eq:suc0}\tag{suc0}
\end{align}
\end{lemma}


\pn From there the ordinal representing 2 is the successor of the ordinal representing 1 is  \\
\( \{ \{ \varnothing \} , \{ \{ \varnothing \} \} \}  \).  
The ordinal representing 3 is 
\( \{ \{ \{ \varnothing \} , \{ \{ \varnothing \} \} \}  , \{ \{ \{ \{ \{ \varnothing \} , \{ \{ \varnothing \} \} \} \} \} \} \).  


\pn The \bemph{ordinal predicate} tests whether a class is an ordinal.  It uses three concepts explained below:  transitive (Tr), well-ordered (We), and the epsilon relation (E).

\begin{metadefinition}\blTitle{df-ord} Define the ordinal predicate, which is
true for a class that is transitive
     and is well-ordered by the epsilon relation. 
\begin{align}
\mm{\vdash} ( {\rm Ord}~ \mc{A} \leftrightarrow ( {\rm Tr}~ \mc{A} \wedge {\rm E}
    ~{\rm We}~ \mc{A} ) )\label{eq:df-ord}\tag{df-ord}
\end{align}
\end{metadefinition}


In this case \bemph{transitive} does not refer to chaining relations, but instead to properties of the successor $suc$ operation.  The union operator $\bigcup$ in effect reduces set containment depth in the opposite way $suc$ increases it.

Suppose I have a set $K$ corresponding to the ordinal number $k$. Its successor will be $\{ K , \{ K \} \}$.  Applying a union operator flattens the set containment depth, as all the elements that themselves are sets, have all their elements combined into a single set.  So union operator on the successor combines the first element, with the members of the set as the second element to get $\{ K , K \}$.   Because sets don't keep repeated elements it becomes $\{ K \}$.\footnote{But that's not the same as $K$!}

\begin{lemma}\blTitle{unisuc} A transitive class is equal to the union of
its successor.  
\begin{align}
\mm{\vdash} \mc{A} \in {\rm V}\quad\Rightarrow\quad \mm{\vdash} ( {\rm Tr}~
    \mc{A} \leftrightarrow \bigcup {\rm suc}~\mc{A} = \mc{A}
    )\label{eq:unisuc}\tag{unisuc}
\end{align}
\end{lemma}

\begin{lemma}\blTitle{suctr} The successor of a transitive class is
transitive.  
\begin{align}
\mm{\vdash} ( {\rm Tr}~\mc{A} \rightarrow {\rm Tr}~{\rm suc}~\mc{A}
    )\label{eq:suctr}\tag{suctr}
\end{align}
\end{lemma}


A set is well ordered by a relation (think $<$) such that every subset of that set has a least element according to that relation.
The theorem below says that $R$ well-orders $A$, if for some subset $B$ of $A$, then there is some element of $B$ without a predecessor.


\begin{lemma}\blTitle{tz6.26} All nonempty (possibly proper) subclasses of
$\mc{A}$, which has a
       well-founded relation $R$, have $R$-minimal elements. 
\begin{align}
\mm{\vdash} ( ( ( R~{\rm We}~\mc{A} \wedge R~{\rm Se}~\mc{A} ) \wedge ( \mc{B}
    \subseteq \mc{A} \wedge \mc{B} \ne \varnothing ) ) \rightarrow \exists
    \msc{y} \in \mc{B}~{\rm Pred} ( R , \mc{B} , \msc{y} ) = \varnothing
    )\label{eq:tz6.26}\tag{tz6.26}
\end{align}
\end{lemma}


Back to ordinal numbers.

\begin{metadefinition}\blTitle{df-on} Define the class of all ordinal numbers. 
\begin{align}
\mm{\vdash} {\rm On} = \{ \msc{x} ~\vert~ {\rm Ord}~ \msc{x}
    \}\label{eq:df-on}\tag{df-on}
\end{align}
\end{metadefinition}


\begin{metadefinition}\blTitle{df-lim} Define the limit ordinal predicate, which
is true for a nonempty ordinal
     that is not a successor (i.e. that is the union of itself).  
\begin{align}
\mm{\vdash} ( {\rm Lim} \mc{A} \leftrightarrow ( {\rm Ord}~ \mc{A} \wedge \mc{A}
    \ne \varnothing \wedge \mc{A} = \bigcup \mc{A} )
    )\label{eq:df-lim}\tag{df-lim}
\end{align}
\end{metadefinition}


\begin{lemma}\blTitle{ordeq} Equality theorem for the ordinal predicate. 
\begin{align}
\mm{\vdash} ( \mc{A} = \mc{B} \rightarrow ( {\rm Ord}~ \mc{A} \leftrightarrow
    {\rm Ord}~ \mc{B} ) )\label{eq:ordeq}\tag{ordeq}
\end{align}
\end{lemma}


\begin{lemma}\blTitle{elong} An ordinal number is an ordinal set. 
\begin{align}
\mm{\vdash} ( \mc{A} \in V \rightarrow ( \mc{A} \in {\rm On} \leftrightarrow
    {\rm Ord}~ \mc{A} ) )\label{eq:elong}\tag{elong}
\end{align}
\end{lemma}


\begin{lemma}\blTitle{eloni} An ordinal number has the ordinal property. 
\begin{align}
\mm{\vdash} ( \mc{A} \in {\rm On} \rightarrow {\rm Ord}~ \mc{A}
    )\label{eq:eloni}\tag{eloni}
\end{align}
\end{lemma}


\begin{lemma}\blTitle{elon2} An ordinal number is an ordinal set. 
\begin{align}
\mm{\vdash} ( \mc{A} \in {\rm On} \leftrightarrow ( {\rm Ord}~ \mc{A} \wedge
    \mc{A} \in {\rm V} ) )\label{eq:elon2}\tag{elon2}
\end{align}
\end{lemma}


\begin{lemma}\blTitle{limeq} Equality theorem for the limit predicate. 
\begin{align}
\mm{\vdash} ( \mc{A} = \mc{B} \rightarrow ( {\rm Lim} \mc{A} \leftrightarrow
    {\rm Lim} \mc{B} ) )\label{eq:limeq}\tag{limeq}
\end{align}
\end{lemma}


\begin{lemma}\blTitle{ordwe} Epsilon well-orders every ordinal. 
\begin{align}
\mm{\vdash} ( {\rm Ord}~ \mc{A} \rightarrow {\rm E} ~{\rm We}~ \mc{A}
    )\label{eq:ordwe}\tag{ordwe}
\end{align}
\end{lemma}


\begin{lemma}\blTitle{ordtr} An ordinal class is transitive.  
\begin{align}
\mm{\vdash} ( {\rm Ord}~ \mc{A} \rightarrow {\rm Tr}~ \mc{A}
    )\label{eq:ordtr}\tag{ordtr}
\end{align}
\end{lemma}


\begin{lemma}\blTitle{ordfr} Epsilon is well-founded on an ordinal class. 
\begin{align}
\mm{\vdash} ( {\rm Ord}~ \mc{A} \rightarrow {\rm E} {\rm Fr} \mc{A}
    )\label{eq:ordfr}\tag{ordfr}
\end{align}
\end{lemma}


\begin{lemma}\blTitle{ordelss} An element of an ordinal class is a subset
of it.  
\begin{align}
\mm{\vdash} ( ( {\rm Ord}~ \mc{A} \wedge \mc{B} \in \mc{A} ) \rightarrow \mc{B}
    \subseteq \mc{A} )\label{eq:ordelss}\tag{ordelss}
\end{align}
\end{lemma}


\begin{lemma}\blTitle{tron} The class of all ordinal numbers is
transitive.  
\begin{align}
\mm{\vdash} {\rm Tr}~ {\rm On}\label{eq:tron}\tag{tron}
\end{align}
\end{lemma}


\begin{lemma}\blTitle{elelsuc} Membership in a successor.  
\begin{align}
\mm{\vdash} ( \mc{A} \in \mc{B} \rightarrow \mc{A} \in {\rm suc}~ \mc{B}
    )\label{eq:elelsuc}\tag{elelsuc}
\end{align}
\end{lemma}




%
%\subsection{Well-Founded Induction}\label{sec:wfinduction}
%
%\begin{theorem}\label{the:tz6.26} All nonempty (possibly proper) subclasses of
%$\mc{A}$, which has a
%       well-founded relation $R$, have $R$-minimal elements.  
%\begin{align}
%\mm{\vdash} ( ( ( R {\rm We} \mc{A} \wedge R {\rm Se} \mc{A} ) \wedge ( \mc{B}
%    \subseteq \mc{A} \wedge \mc{B} \ne \varnothing ) ) \rightarrow \exists
%    \msc{y} \in \mc{B} {\rm Pred} ( R , \mc{B} , \msc{y} ) = \varnothing
%    )\label{eq:tz6.26}\tag{tz6.26}
%\end{align}
%\end{theorem}
%
%\begin{theorem}\label{the:wfi} The Principle of Well-Founded Induction. 
% This principle states that if $\mc{B}$ is a
%       subclass of a well-ordered class $\mc{A}$ with the property that every
%       element of $\mc{B}$ whose inital segment is included in $\mc{A}$ is itself
%       equal to $\mc{A}$.
%\begin{align}
%\mm{\vdash} ( ( ( R {\rm We} \mc{A} \wedge R {\rm Se} \mc{A} ) \wedge ( \mc{B}
%    \subseteq \mc{A} \wedge \forall \msc{y} \in \mc{A} ( {\rm Pred} ( R ,
%    \mc{A} , \msc{y} ) \subseteq \mc{B} \rightarrow \msc{y} \in \mc{B} ) ) )
%    \rightarrow \mc{A} = \mc{B} )\label{eq:wfi}\tag{wfi}
%\end{align}
%\end{theorem}
%
%\begin{theorem}\label{the:wfisg} Well-Founded Induction Schema.  If a property
%passes from all elements
%       less than $\msc{y}$ of a well-founded class $\mc{A}$ to $\msc{y}$
%itself (induction
%       hypothesis), then the property holds for all elements of $\mc{A}$.
%       (Contributed by Scott Fenton, 11-Feb-2011.)
%\begin{align}
%\mm{\vdash} ( \msc{y} \in \mc{A} \rightarrow ( \forall \msc{z} \in {\rm Pred} (
%    R , \mc{A} , \msc{y} ) {[} \msc{z} / \msc{y}
%    {]} \mw{\varphi} \rightarrow \mw{\varphi} )
%    )\quad\Rightarrow\quad \mm{\vdash} ( ( R {\rm We} \mc{A} \wedge R {\rm Se}
%    \mc{A} ) \rightarrow \forall \msc{y} \in \mc{A} \mw{\varphi}
%    )\label{eq:wfisg}\tag{wfisg}
%\end{align}
%\end{theorem}
%
%\subsection{Russell's definition description binder (inverted iota)}\label{sec:russell}
%
%The one-and-only operator: $\mathrm{\rotatebox[origin=C]{180}{$\iota$}} $.
%
%% This LaTeX file was created by Metamath on 9-Oct-2021 11:49 AM.
\begin{metadefinition}\blTitle{cio} Extend class notation with Russell's
definition description binder
     (inverted iota).
\begin{align}
{\rm \mm{class}~} ( \mathrm{\rotatebox[origin=C]{180}{$\iota$}} \msc{x}
    \mw{\varphi} )\label{eq:cio}\tag{cio}
\end{align}
\end{metadefinition}

\begin{lemma}\blTitle{iotajust} Soundness justification theorem for {\tt
df-iota}.  (Contributed by Andrew
       Salmon, 29-Jun-2011.)
\begin{align}
\mm{\vdash} \bigcup \{ \msc{y} ~\vert~ \{ \msc{x} ~\vert~ \mw{\varphi} \} = \{
    \msc{y} \} \} = \bigcup \{ \msc{z} ~\vert~ \{ \msc{x} ~\vert~ \mw{\varphi}
    \} = \{ \msc{z} \} \}\label{eq:iotajust}\tag{iotajust}
\end{align}
\end{lemma}


\begin{metadefinition}\blTitle{df-iota} Define Russell's definition description
binder, which can be read as
       ``the unique $\msc{x}$ such that $\mw{\varphi}$," where $\mw{\varphi}$
ordinarily contains
       $\msc{x}$ as a free variable.  Our definition is meaningful only when
there
       is exactly one $\msc{x}$ such that $\mw{\varphi}$ is true (see {\tt
iotaval});
       otherwise, it evaluates to the empty set (see {\tt iotanul}).  Russell
used
       the inverted iota symbol $\mathrm{\rotatebox[origin=C]{180}{$\iota$}}$
to represent the binder.

       Sometimes proofs need to expand an iota-based definition.  That is,
       given ``X = the x for which ... x ... x ..." holds, the proof needs to
       get to ``...  X ...  X ...".  A general strategy to do this is to use
       {\tt riotacl2} (or {\tt iotacl} for unbounded iota), as demonstrated in
the
       proof of {\tt supub}.  This can be easier than applying {\tt riotasbc}
or a
       version that applies an explicit substitution, because substituting an
       iota into its own property always has a bound variable clash which must
       be first renamed or else guarded with NF.

       (Contributed by Andrew Salmon, 30-Jun-2011.)
\begin{align}
\mm{\vdash} ( \mathrm{\rotatebox[origin=C]{180}{$\iota$}} \msc{x} \mw{\varphi}
    ) = \bigcup \{ \msc{y} ~\vert~ \{ \msc{x} ~\vert~ \mw{\varphi} \} = \{
    \msc{y} \} \}\label{eq:df-iota}\tag{df-iota}
\end{align}
\end{metadefinition}

\begin{lemma}\blTitle{dfiota2} Alternate definition for descriptions. 
Definition 8.18 in [Quine]
       p. 56.  (Contributed by Andrew Salmon, 30-Jun-2011.)
\begin{align}
\mm{\vdash} ( \mathrm{\rotatebox[origin=C]{180}{$\iota$}} \msc{x} \mw{\varphi}
    ) = \bigcup \{ \msc{y} ~\vert~ \forall \msc{x} ( \mw{\varphi}
    \leftrightarrow \msc{x} = \msc{y} ) \}\label{eq:dfiota2}\tag{dfiota2}
\end{align}
\end{lemma}


\begin{lemma}\blTitle{nfiota1} Bound-variable hypothesis builder for the
$\mathrm{\rotatebox[origin=C]{180}{$\iota$}}$ class.  (Contributed
       by Andrew Salmon, 11-Jul-2011.)  (Revised by Mario Carneiro,
       15-Oct-2016.)
\begin{align}
\mm{\vdash} % \underline{\Finv} 
\msc{x} (
    \mathrm{\rotatebox[origin=C]{180}{$\iota$}} \msc{x} \mw{\varphi}
    )\label{eq:nfiota1}\tag{nfiota1}
\end{align}
\end{lemma}


\begin{lemma}\blTitle{nfiotad} Deduction version of {\tt nfiota}. 
(Contributed by NM, 18-Feb-2013.)
\begin{align}
\mm{\vdash} \Finv 
\msc{y} \mw{\varphi}\quad\&\quad \mm{\vdash} ( \mw{\varphi}
    \rightarrow \Finv 
    \msc{x} \mw{\psi} )\quad\Rightarrow\quad \mm{\vdash} (
    \mw{\varphi} \rightarrow % \underline{\Finv} 
    \msc{x} (
    \mathrm{\rotatebox[origin=C]{180}{$\iota$}} \msc{y} \mw{\psi} )
    )\label{eq:nfiotad}\tag{nfiotad}
\end{align}
\end{lemma}


\begin{lemma}\blTitle{nfiota} Bound-variable hypothesis builder for the
$\mathrm{\rotatebox[origin=C]{180}{$\iota$}}$ class.  (Contributed
       by NM, 23-Aug-2011.)
\begin{align}
\mm{\vdash} % \Finv 
\msc{x} \mw{\varphi}\quad\Rightarrow\quad \mm{\vdash}
  %  \underline{\Finv} 
  \msc{x} ( \mathrm{\rotatebox[origin=C]{180}{$\iota$}}
    \msc{y} \mw{\varphi} )\label{eq:nfiota}\tag{nfiota}
\end{align}
\end{lemma}


\begin{lemma}\blTitle{iotajust} Soundness justification theorem for {\tt
df-iota}.  (Contributed by Andrew
       Salmon, 29-Jun-2011.)
\begin{align}
\mm{\vdash} \bigcup \{ \msc{y} ~\vert~ \{ \msc{x} ~\vert~ \mw{\varphi} \} = \{
    \msc{y} \} \} = \bigcup \{ \msc{z} ~\vert~ \{ \msc{x} ~\vert~ \mw{\varphi}
    \} = \{ \msc{z} \} \}\label{eq:iotajust}\tag{iotajust}
\end{align}
\end{lemma}


\begin{lemma}\blTitle{iotaeq} Equality theorem for descriptions. 
(Contributed by Andrew Salmon,
       30-Jun-2011.)
\begin{align}
\mm{\vdash} ( \forall \msc{x} \msc{x} = \msc{y} \rightarrow (
    \mathrm{\rotatebox[origin=C]{180}{$\iota$}} \msc{x} \mw{\varphi} ) = (
    \mathrm{\rotatebox[origin=C]{180}{$\iota$}} \msc{y} \mw{\varphi} )
    )\label{eq:iotaeq}\tag{iotaeq}
\end{align}
\end{lemma}


%\begin{lemma}\blTitle{iotabi} Equivalence theorem for descriptions. 
%(Contributed by Andrew Salmon,
%       30-Jun-2011.)
%\begin{align}
%\mm{\vdash} ( \forall \msc{x} ( \mw{\varphi} \leftrightarrow \mw{\psi} )
%    \rightarrow ( \mathrm{\rotatebox[origin=C]{180}{$\iota$}} \msc{x}
%    \mw{\varphi} ) = ( \mathrm{\rotatebox[origin=C]{180}{$\iota$}} \msc{x}
%    \mw{\psi} ) )\label{eq:iotabi}\tag{iotabi}
%\end{align}
%\end{lemma}

%
%\begin{lemma}\blTitle{iotaval} Theorem 8.19 in [Quine] p. 57.  This
%theorem is the fundamental property
%       of iota.  (Contributed by Andrew Salmon, 11-Jul-2011.)
%\begin{align}
%\mm{\vdash} ( \forall \msc{x} ( \mw{\varphi} \leftrightarrow \msc{x} = \msc{y}
%    ) \rightarrow ( \mathrm{\rotatebox[origin=C]{180}{$\iota$}} \msc{x}
%    \mw{\varphi} ) = \msc{y} )\label{eq:iotaval}\tag{iotaval}
%\end{align}
%\end{lemma}

%
%\begin{lemma}\blTitle{iotauni} Equivalence between two different forms of
%$\mathrm{\rotatebox[origin=C]{180}{$\iota$}}$.  (Contributed by
%       Andrew Salmon, 12-Jul-2011.)
%\begin{align}
%\mm{\vdash} ( \exists{!} \msc{x} \mw{\varphi} \rightarrow (
%    \mathrm{\rotatebox[origin=C]{180}{$\iota$}} \msc{x} \mw{\varphi} ) =
%    \bigcup \{ \msc{x} ~\vert~ \mw{\varphi} \} )\label{eq:iotauni}\tag{iotauni}
%\end{align}
%\end{lemma}

%
%\begin{lemma}\blTitle{iotaint} Equivalence between two different forms of
%$\mathrm{\rotatebox[origin=C]{180}{$\iota$}}$.  (Contributed by
%       Mario Carneiro, 24-Dec-2016.)
%\begin{align}
%\mm{\vdash} ( \exists{!} \msc{x} \mw{\varphi} \rightarrow (
%    \mathrm{\rotatebox[origin=C]{180}{$\iota$}} \msc{x} \mw{\varphi} ) =
%    \bigcap \{ \msc{x} ~\vert~ \mw{\varphi} \} )\label{eq:iotaint}\tag{iotaint}
%\end{align}
%\end{lemma}

%
%\begin{lemma}\blTitle{iota1} Property of iota.  (Contributed by NM,
%23-Aug-2011.)  (Revised by Mario
%       Carneiro, 23-Dec-2016.)
%\begin{align}
%\mm{\vdash} ( \exists{!} \msc{x} \mw{\varphi} \rightarrow ( \mw{\varphi}
%    \leftrightarrow ( \mathrm{\rotatebox[origin=C]{180}{$\iota$}} \msc{x}
%    \mw{\varphi} ) = \msc{x} ) )\label{eq:iota1}\tag{iota1}
%\end{align}
%\end{lemma}

%
%\begin{lemma}\blTitle{iotanul} Theorem 8.22 in [Quine] p. 57.  This
%theorem is the result if there
%       isn't exactly one $\msc{x}$ that satisfies $\mw{\varphi}$. 
%(Contributed by Andrew
%       Salmon, 11-Jul-2011.)
%\begin{align}
%\mm{\vdash} ( \lnot \exists{!} \msc{x} \mw{\varphi} \rightarrow (
%    \mathrm{\rotatebox[origin=C]{180}{$\iota$}} \msc{x} \mw{\varphi} ) =
%    \varnothing )\label{eq:iotanul}\tag{iotanul}
%\end{align}
%\end{lemma}

%
%\begin{lemma}\blTitle{iotassuni} The
%$\mathrm{\rotatebox[origin=C]{180}{$\iota$}}$ class is a subset of the union
%of all elements satisfying
%       $\mw{\varphi}$.  (Contributed by Mario Carneiro, 24-Dec-2016.)
%\begin{align}
%\mm{\vdash} ( \mathrm{\rotatebox[origin=C]{180}{$\iota$}} \msc{x} \mw{\varphi}
%    ) \subseteq \bigcup \{ \msc{x} ~\vert~ \mw{\varphi}
%    \}\label{eq:iotassuni}\tag{iotassuni}
%\end{align}
%\end{lemma}

%
%\begin{lemma}\blTitle{iotaex} Theorem 8.23 in [Quine] p. 58.  This theorem
%proves the existence of the
%       $\mathrm{\rotatebox[origin=C]{180}{$\iota$}}$ class under our
%definition.  (Contributed by Andrew Salmon,
%       11-Jul-2011.)
%\begin{align}
%\mm{\vdash} ( \mathrm{\rotatebox[origin=C]{180}{$\iota$}} \msc{x} \mw{\varphi}
%    ) \in {\rm V}\label{eq:iotaex}\tag{iotaex}
%\end{align}
%\end{lemma}

%
%\begin{lemma}\blTitle{iota4} Theorem *14.22 in [WhiteheadRussell] p. 190. 
%(Contributed by Andrew
%       Salmon, 12-Jul-2011.)
%\begin{align}
%\mm{\vdash} ( \exists{!} \msc{x} \mw{\varphi} \rightarrow \underaccent{\dot}
%  {[}
%    ( \mathrm{\rotatebox[origin=C]{180}{$\iota$}} \msc{x} \mw{\varphi} ) /
%    \msc{x} \underaccent{\dot}
%    {]} \mw{\varphi} )\label{eq:iota4}\tag{iota4}
%\end{align}
%\end{lemma}

%
%\begin{lemma}\blTitle{iota4an} Theorem *14.23 in [WhiteheadRussell] p.
%191.  (Contributed by Andrew
%     Salmon, 12-Jul-2011.)
%\begin{align}
%\mm{\vdash} ( \exists{!} \msc{x} ( \mw{\varphi} \wedge \mw{\psi} ) \rightarrow
%   \underaccent{\dot}
%   {[} ( \mathrm{\rotatebox[origin=C]{180}{$\iota$}} \msc{x}
%    ( \mw{\varphi} \wedge \mw{\psi} ) ) / \msc{x} \underaccent{\dot}
%    {]}
%    \mw{\varphi} )\label{eq:iota4an}\tag{iota4an}
%\end{align}
%\end{lemma}

%
%\begin{lemma}\blTitle{iota5} A method for computing iota.  (Contributed by
%NM, 17-Sep-2013.)
%\begin{align}
%\mm{\vdash} ( ( \mw{\varphi} \wedge \mc{A} \in V ) \rightarrow ( \mw{\psi}
%    \leftrightarrow \msc{x} = \mc{A} ) )\quad\Rightarrow\quad \mm{\vdash} ( (
%    \mw{\varphi} \wedge \mc{A} \in V ) \rightarrow (
%    \mathrm{\rotatebox[origin=C]{180}{$\iota$}} \msc{x} \mw{\psi} ) = \mc{A}
%    )\label{eq:iota5}\tag{iota5}
%\end{align}
%\end{lemma}

%
%\begin{lemma}\blTitle{iotabidv} Formula-building deduction rule for iota. 
%(Contributed by NM,
%       20-Aug-2011.)
%\begin{align}
%\mm{\vdash} ( \mw{\varphi} \rightarrow ( \mw{\psi} \leftrightarrow \mw{\chi} )
%    )\quad\Rightarrow\quad \mm{\vdash} ( \mw{\varphi} \rightarrow (
%    \mathrm{\rotatebox[origin=C]{180}{$\iota$}} \msc{x} \mw{\psi} ) = (
%    \mathrm{\rotatebox[origin=C]{180}{$\iota$}} \msc{x} \mw{\chi} )
%    )\label{eq:iotabidv}\tag{iotabidv}
%\end{align}
%\end{lemma}

%
%\begin{lemma}\blTitle{iotabii} Formula-building deduction rule for iota. 
%(Contributed by Mario
%       Carneiro, 2-Oct-2015.)
%\begin{align}
%\mm{\vdash} ( \mw{\varphi} \leftrightarrow \mw{\psi} )\quad\Rightarrow\quad
%    \mm{\vdash} ( \mathrm{\rotatebox[origin=C]{180}{$\iota$}} \msc{x}
%    \mw{\varphi} ) = ( \mathrm{\rotatebox[origin=C]{180}{$\iota$}} \msc{x}
%    \mw{\psi} )\label{eq:iotabii}\tag{iotabii}
%\end{align}
%\end{lemma}

%
%\begin{lemma}\blTitle{iotacl} Membership law for descriptions.
%
%     This can useful for expanding an unbounded iota-based definition (see
%     {\tt df-iota}).  If you have a bounded iota-based definition, {\tt
%riotacl2} may
%     be useful.
%
%     (Contributed by Andrew Salmon, 1-Aug-2011.)
%\begin{align}
%\mm{\vdash} ( \exists{!} \msc{x} \mw{\varphi} \rightarrow (
%    \mathrm{\rotatebox[origin=C]{180}{$\iota$}} \msc{x} \mw{\varphi} ) \in \{
%    \msc{x} ~\vert~ \mw{\varphi} \} )\label{eq:iotacl}\tag{iotacl}
%\end{align}
%\end{lemma}

%
%\begin{lemma}\blTitle{iota2df} A condition that allows us to represent
%``the unique element such that
%         $\mw{\varphi}$" with a class expression $\mc{A}$.  (Contributed by
%NM,
%         30-Dec-2014.)
%\begin{align}
%\mm{\vdash} ( \mw{\varphi} \rightarrow \mc{B} \in V )\quad\&\quad \mm{\vdash} (
%    \mw{\varphi} \rightarrow \exists{!} \msc{x} \mw{\psi} )\quad\&\quad
%    \mm{\vdash} ( ( \mw{\varphi} \wedge \msc{x} = \mc{B} ) \rightarrow (
%    \mw{\psi} \leftrightarrow \mw{\chi} ) ) \\
%    \quad\&\quad \mm{\vdash}
%     \Finv 
%    \msc{x} \mw{\varphi}\quad\&\quad \mm{\vdash} ( \mw{\varphi} \rightarrow
%    \Finv  
%    \msc{x} \mw{\chi} )\quad\&\quad \mm{\vdash} ( \mw{\varphi}
%    \rightarrow 
%    \underline{\Finv} 
%    \msc{x} \mc{B} )\quad\Rightarrow\quad
%    \mm{\vdash} ( \mw{\varphi} \rightarrow ( \mw{\chi} \leftrightarrow (
%    \mathrm{\rotatebox[origin=C]{180}{$\iota$}} \msc{x} \mw{\psi} ) = \mc{B} )
%    )\label{eq:iota2df}\tag{iota2df}
%\end{align}
%\end{lemma}

%
%\begin{lemma}\blTitle{iota2d} A condition that allows us to represent
%``the unique element such that
%       $\mw{\varphi}$" with a class expression $\mc{A}$.  (Contributed by NM,
%       30-Dec-2014.)
%\begin{align}
%\mm{\vdash} ( \mw{\varphi} \rightarrow \mc{B} \in V )\quad\&\quad \mm{\vdash} (
%    \mw{\varphi} \rightarrow \exists{!} \msc{x} \mw{\psi} )\quad\&\quad
%    \mm{\vdash} ( ( \mw{\varphi} \wedge \msc{x} = \mc{B} ) \rightarrow (
%    \mw{\psi} \leftrightarrow \mw{\chi} ) )\quad\Rightarrow\quad \mm{\vdash} (
%    \mw{\varphi} \rightarrow ( \mw{\chi} \leftrightarrow (
%    \mathrm{\rotatebox[origin=C]{180}{$\iota$}} \msc{x} \mw{\psi} ) = \mc{B} )
%    )\label{eq:iota2d}\tag{iota2d}
%\end{align}
%\end{lemma}

%
%\begin{lemma}\blTitle{iota2} The unique element such that $\mw{\varphi}$. 
%(Contributed by Jeff Madsen,
%       1-Jun-2011.)  (Revised by Mario Carneiro, 23-Dec-2016.)
%\begin{align}
%\mm{\vdash} ( \msc{x} = \mc{A} \rightarrow ( \mw{\varphi} \leftrightarrow
%    \mw{\psi} ) )\quad\Rightarrow\quad \mm{\vdash} ( ( \mc{A} \in \mc{B} \wedge
%    \exists{!} \msc{x} \mw{\varphi} ) \rightarrow ( \mw{\psi} \leftrightarrow (
%    \mathrm{\rotatebox[origin=C]{180}{$\iota$}} \msc{x} \mw{\varphi} ) = \mc{A}
%    ) )\label{eq:iota2}\tag{iota2}
%\end{align}
%\end{lemma}

%
%\begin{lemma}\blTitle{iota5f} A method for computing iota.  (Contributed
%by Scott Fenton,
%       13-Dec-2017.)
%\begin{align}
%\mm{\vdash} \Finv  
%\msc{x} \mw{\varphi}\quad\&\quad \mm{\vdash}
%   \underline{\Finv} 
%  \msc{x} \mc{A}\quad\&\quad \mm{\vdash} ( ( \mw{\varphi}
%    \wedge \mc{A} \in V ) \rightarrow ( \mw{\psi} \leftrightarrow \msc{x} =
%    \mc{A} ) )\quad\Rightarrow\quad \mm{\vdash} ( ( \mw{\varphi} \wedge \mc{A}
%    \in V ) \rightarrow ( \mathrm{\rotatebox[origin=C]{180}{$\iota$}} \msc{x}
%    \mw{\psi} ) = \mc{A} )\label{eq:iota5f}\tag{iota5f}
%\end{align}
%\end{lemma}

%
%\begin{lemma}\blTitle{iotain} Equivalence between two different forms of
%$\mathrm{\rotatebox[origin=C]{180}{$\iota$}}$.  (Contributed by
%       Andrew Salmon, 15-Jul-2011.)
%\begin{align}
%\mm{\vdash} ( \exists{!} \msc{x} \mw{\varphi} \rightarrow \bigcap \{ \msc{x}
%    ~\vert~ \mw{\varphi} \} = ( \mathrm{\rotatebox[origin=C]{180}{$\iota$}}
%    \msc{x} \mw{\varphi} ) )\label{eq:iotain}\tag{iotain}
%\end{align}
%\end{lemma}

%
%\begin{lemma}\blTitle{iotaexeu} The iota class exists.  This theorem does
%not require {\tt ax-nul} for its
%       proof.  (Contributed by Andrew Salmon, 11-Jul-2011.)
%\begin{align}
%\mm{\vdash} ( \exists{!} \msc{x} \mw{\varphi} \rightarrow (
%    \mathrm{\rotatebox[origin=C]{180}{$\iota$}} \msc{x} \mw{\varphi} ) \in {\rm
%    V} )\label{eq:iotaexeu}\tag{iotaexeu}
%\end{align}
%\end{lemma}

%
%\begin{lemma}\blTitle{iotasbc} Definition *14.01 in [WhiteheadRussell] p.
%184.  In Principia
%       Mathematica, Russell and Whitehead define
%$\mathrm{\rotatebox[origin=C]{180}{$\iota$}}$ in terms of a
%       function of $( \mathrm{\rotatebox[origin=C]{180}{$\iota$}} \msc{x}
%\mw{\varphi} )$.  Their definition differs in that a
%       function of $( \mathrm{\rotatebox[origin=C]{180}{$\iota$}} \msc{x}
%\mw{\varphi} )$ evaluates to ``false" when there isn't a
%       single $\msc{x}$ that satisfies $\mw{\varphi}$.  (Contributed by Andrew
%Salmon,
%       11-Jul-2011.)
%\begin{align}
%\mm{\vdash} ( \exists{!} \msc{x} \mw{\varphi} \rightarrow (
%    \underaccent{\dot}
%    {[} ( \mathrm{\rotatebox[origin=C]{180}{$\iota$}} \msc{x}
%    \mw{\varphi} ) / \msc{y} \underaccent{\dot}
%    {]} \mw{\psi} \leftrightarrow
%    \exists \msc{y} ( \forall \msc{x} ( \mw{\varphi} \leftrightarrow \msc{x} =
%    \msc{y} ) \wedge \mw{\psi} ) ) )\label{eq:iotasbc}\tag{iotasbc}
%\end{align}
%\end{lemma}

%
%\begin{lemma}\blTitle{iotasbc2} Theorem *14.111 in [WhiteheadRussell] p.
%184.  (Contributed by Andrew
%       Salmon, 11-Jul-2011.)
%\begin{align}
%\mm{\vdash} ( ( \exists{!} \msc{x} \mw{\varphi} \wedge \exists{!} \msc{x}
%    \mw{\psi} ) \rightarrow ( \underaccent{\dot}{[} (
%    \mathrm{\rotatebox[origin=C]{180}{$\iota$}} \msc{x} \mw{\varphi} ) /
%    \msc{y} \underaccent{\dot}
%    {]} \underaccent{\dot}{[} (
%    \mathrm{\rotatebox[origin=C]{180}{$\iota$}} \msc{x} \mw{\psi} ) / \msc{z}
%    \underaccent{\dot}{]} \mw{\chi} \leftrightarrow \exists \msc{y} \exists
%    \msc{z} ( \forall \msc{x} ( \mw{\varphi} \leftrightarrow \msc{x} = \msc{y}
%    ) \wedge \forall \msc{x} ( \mw{\psi} \leftrightarrow \msc{x} = \msc{z} )
%    \wedge \mw{\chi} ) ) )\label{eq:iotasbc2}\tag{iotasbc2}
%\end{align}
%\end{lemma}

%
%\begin{lemma}\blTitle{iotaequ} Theorem *14.2 in [WhiteheadRussell] p. 189.
%(Contributed by Andrew
%       Salmon, 11-Jul-2011.)
%\begin{align}
%\mm{\vdash} ( \mathrm{\rotatebox[origin=C]{180}{$\iota$}} \msc{x} \msc{x} =
%    \msc{y} ) = \msc{y}\label{eq:iotaequ}\tag{iotaequ}
%\end{align}
%\end{lemma}

%
%\begin{lemma}\blTitle{iotavalb} Theorem *14.202 in [WhiteheadRussell] p.
%189.  A biconditional version
%       of {\tt iotaval}.  (Contributed by Andrew Salmon, 11-Jul-2011.)
%\begin{align}
%\mm{\vdash} ( \exists{!} \msc{x} \mw{\varphi} \rightarrow ( \forall \msc{x} (
%    \mw{\varphi} \leftrightarrow \msc{x} = \msc{y} ) \leftrightarrow (
%    \mathrm{\rotatebox[origin=C]{180}{$\iota$}} \msc{x} \mw{\varphi} ) =
%    \msc{y} ) )\label{eq:iotavalb}\tag{iotavalb}
%\end{align}
%\end{lemma}

%
%\begin{lemma}\blTitle{iotasbc5} Theorem *14.205 in [WhiteheadRussell] p.
%190.  (Contributed by Andrew
%       Salmon, 11-Jul-2011.)
%\begin{align}
%\mm{\vdash} ( \exists{!} \msc{x} \mw{\varphi} \rightarrow (
%   \underaccent{\dot}
%   {[} ( \mathrm{\rotatebox[origin=C]{180}{$\iota$}} \msc{x}
%    \mw{\varphi} ) / \msc{y} \underaccent{\dot}
%    {]} \mw{\psi} \leftrightarrow
%    \exists \msc{y} ( \msc{y} = ( \mathrm{\rotatebox[origin=C]{180}{$\iota$}}
%    \msc{x} \mw{\varphi} ) \wedge \mw{\psi} ) )
%    )\label{eq:iotasbc5}\tag{iotasbc5}
%\end{align}
%\end{lemma}

%
%\begin{lemma}\blTitle{iotavalsb} Theorem *14.242 in [WhiteheadRussell] p.
%192.  (Contributed by Andrew
%       Salmon, 11-Jul-2011.)
%\begin{align}
%\mm{\vdash} ( \forall \msc{x} ( \mw{\varphi} \leftrightarrow \msc{x} = \msc{y}
%    ) \rightarrow ( \underaccent{\dot}
%    {[} \msc{y} / \msc{z}
%    \underaccent{\dot}
%    {]} \mw{\psi} \leftrightarrow \underaccent{\dot}
%    {[} (
%    \mathrm{\rotatebox[origin=C]{180}{$\iota$}} \msc{x} \mw{\varphi} ) /
%    \msc{z} \underaccent{\dot}
%    {]} \mw{\psi} )
%    )\label{eq:iotavalsb}\tag{iotavalsb}
%\end{align}
%\end{lemma}

%
%\begin{lemma}\blTitle{iotasbcq} Theorem *14.272 in [WhiteheadRussell] p.
%193.  (Contributed by Andrew
%       Salmon, 11-Jul-2011.)
%\begin{align}
%\mm{\vdash} ( \forall \msc{x} ( \mw{\varphi} \leftrightarrow \mw{\psi} )
%    \rightarrow ( \underaccent{\dot}
%    {[} (
%    \mathrm{\rotatebox[origin=C]{180}{$\iota$}} \msc{x} \mw{\varphi} ) /
%    \msc{y} \underaccent{\dot}
%    {]} \mw{\chi} \leftrightarrow
%    \underaccent{\dot}
%    {[} ( \mathrm{\rotatebox[origin=C]{180}{$\iota$}} \msc{x}
%    \mw{\psi} ) / \msc{y} \underaccent{\dot}
%    {]} \mw{\chi} )
%    )\label{eq:iotasbcq}\tag{iotasbcq}
%\end{align}
%\end{lemma}

%


%
%\subsection{Recursive Definition Generator}\label{sec:recdefgen}
%
%% This LaTeX file was created by Metamath on 9-Oct-2021 1:49 PM.

\begin{metadefinition}\blTitle{crdg} Extend class notation with the recursive
definition generator, with
     characteristic function $F$ and initial value $I$.
\begin{align}
{\rm \mm{class}~} {\rm rec} ( F , I )\label{eq:crdg}\tag{crdg}
\end{align}
\end{metadefinition}

\begin{metadefinition}\blTitle{df-rdg} Define a recursive definition generator
on ${\rm On}$ (the class of ordinal
       numbers) with characteristic function $F$ and initial value $I$.
       This combines functions $F$ in {\tt tfr1} and $G$ in {\tt tz7.44-1}
into one
       definition.  This rather amazing operation allows us to define, with
       compact direct definitions, functions that are usually defined in
       textbooks only with indirect self-referencing recursive definitions.  A
       recursive definition requires advanced metalogic to justify - in
       particular, eliminating a recursive definition is very difficult and
       often not even shown in textbooks.  On the other hand, the elimination
       of a direct definition is a matter of simple mechanical substitution.
       The price paid is the daunting complexity of our ${\rm rec}$ operation
       (especially when {\tt df-recs} that it is built on is also eliminated).
But
       once we get past this hurdle, definitions that would otherwise be
       recursive become relatively simple, as in for example {\tt oav}, from
which
       we prove the recursive textbook definition as theorems {\tt oa0}, {\tt
oasuc},
       and {\tt oalim} (with the help of theorems {\tt rdg0}, {\tt rdgsuc},
and
       {\tt rdglim2a}).  We can also restrict the ${\rm rec}$ operation to
define
       otherwise recursive functions on the natural numbers ${\rm \omega}$;
see
       {\tt fr0g} and {\tt frsuc}.  Our ${\rm rec}$ operation apparently does
not appear
       in published literature, although closely related is Definition 25.2 of
       [Quine] p. 177, which he uses to ``turn...a recursion into a genuine or
       direct definition" (p. 174).  Note that the ${\rm if}$ operations (see
       {\tt df-if}) select cases based on whether the domain of $g$ is zero, a
       successor, or a limit ordinal.

       An important use of this definition is in the recursive sequence
       generator {\tt df-seq} on the natural numbers (as a subset of the
complex
       numbers), allowing us to define, with direct definitions, recursive
       infinite sequences such as the factorial function {\tt df-fac} and
integer
       powers {\tt df-exp}.

       -Note:  We introduce- ${\rm rec}$-with the philosophical goal of being-
       -able to eliminate all definitions with direct mechanical substitution-
       -and to verify easily the soundness of definitions.  Metamath itself-
       -has no built-in technical limitation that prevents multiple-part-
       -recursive definitions in the traditional textbook style-. 
(Contributed
       by NM, 9-Apr-1995.)  (Revised by Mario Carneiro, 9-May-2015.)
\begin{align}
\mm{\vdash} {\rm rec} ( F , I ) = \mathrm{recs} ( ( g \in {\rm V} \mapsto {\rm
    if} ( g = \varnothing , I , {\rm if} ( {\rm Lim}~ {\rm dom}~ g , \bigcup {\rm
    ran} g , ( F ` ( g ` \bigcup {\rm dom}~ g ) ) ) ) )
    )\label{eq:df-rdg}\tag{df-rdg}
\end{align}
\end{metadefinition}

\begin{lemma}\blTitle{rdgeq1} Equality theorem for the recursive
definition generator.  (Contributed
       by NM, 9-Apr-1995.)  (Revised by Mario Carneiro, 9-May-2015.)
\begin{align}
\mm{\vdash} ( F = G \rightarrow {\rm rec} ( F , \mc{A} ) = {\rm rec} ( G ,
    \mc{A} ) )\label{eq:rdgeq1}\tag{rdgeq1}
\end{align}
\end{lemma}


\begin{lemma}\blTitle{rdgeq1} Equality theorem for the recursive
definition generator.  (Contributed
       by NM, 9-Apr-1995.)  (Revised by Mario Carneiro, 9-May-2015.)
\begin{align}
\mm{\vdash} ( F = G \rightarrow {\rm rec} ( F , \mc{A} ) = {\rm rec} ( G ,
    \mc{A} ) )\label{eq:rdgeq1}\tag{rdgeq1}
\end{align}
\end{lemma}


\begin{lemma}\blTitle{rdgeq2} Equality theorem for the recursive
definition generator.  (Contributed
       by NM, 9-Apr-1995.)  (Revised by Mario Carneiro, 9-May-2015.)
\begin{align}
\mm{\vdash} ( \mc{A} = \mc{B} \rightarrow {\rm rec} ( F , \mc{A} ) = {\rm rec}
    ( F , \mc{B} ) )\label{eq:rdgeq2}\tag{rdgeq2}
\end{align}
\end{lemma}


\begin{lemma}\blTitle{rdgeq12} Equality theorem for the recursive
definition generator.  (Contributed by
     Scott Fenton, 28-Apr-2012.)
\begin{align}
\mm{\vdash} ( ( F = G \wedge \mc{A} = \mc{B} ) \rightarrow {\rm rec} ( F ,
    \mc{A} ) = {\rm rec} ( G , \mc{B} ) )\label{eq:rdgeq12}\tag{rdgeq12}
\end{align}
\end{lemma}


\begin{lemma}\blTitle{rdglem1} Lemma used with the recursive definition
generator.  This is a trivial
       lemma that just changes bound variables for later use.  (Contributed by
       NM, 9-Apr-1995.)
\begin{align}
\mm{\vdash} \{ f ~\vert~ \exists \msc{x} \in {\rm On} ( f {\rm Fn} \msc{x}
    \wedge \forall \msc{y} \in \msc{x} ( f ` \msc{y} ) = 
    ( G ` ( f \restriction
    \msc{y} ) ) ) \} = \\
     \{ g ~\vert~ \exists \msc{z} \in {\rm On} ( g {\rm Fn}
    \msc{z} \wedge \forall \msc{w} \in \msc{z} ( g ` \msc{w} ) = ( G ` ( g
    \restriction \msc{w} ) ) ) \}\label{eq:rdglem1}\tag{rdglem1}
\end{align}
\end{lemma}


\begin{lemma}\blTitle{rdgfun} The recursive definition generator is a
function.  (Contributed by Mario
       Carneiro, 16-Nov-2014.)
\begin{align}
\mm{\vdash} {\rm Fun} {\rm rec} ( F , \mc{A} )\label{eq:rdgfun}\tag{rdgfun}
\end{align}
\end{lemma}


\begin{lemma}\blTitle{rdgdmlim} The domain of the recursive definition
generator is a limit ordinal.
       (Contributed by NM, 16-Nov-2014.)
\begin{align}
\mm{\vdash} {\rm Lim}~ {\rm dom}~ {\rm rec} ( F , \mc{A}
    )\label{eq:rdgdmlim}\tag{rdgdmlim}
\end{align}
\end{lemma}


\begin{lemma}\blTitle{rdgfnon} The recursive definition generator is a
function on ordinal numbers.
       (Contributed by NM, 9-Apr-1995.)  (Revised by Mario Carneiro,
       9-May-2015.)
\begin{align}
\mm{\vdash} {\rm rec} ( F , \mc{A} ) {\rm Fn} {\rm
    On}\label{eq:rdgfnon}\tag{rdgfnon}
\end{align}
\end{lemma}


\begin{lemma}\blTitle{rdgvalg} Value of the recursive definition
generator.  (Contributed by NM,
       9-Apr-1995.)  (Revised by Mario Carneiro, 8-Sep-2013.)
\begin{align}
\mm{\vdash} ( \mc{B} \in {\rm dom}~ {\rm rec} ( F , \mc{A} ) \rightarrow ( {\rm
    rec} ( F , \mc{A} ) ` \mc{B} ) = ( ( g \in {\rm V} \mapsto {\rm if} ( g =
    \varnothing , \mc{A} , {\rm if} ( {\rm Lim}~ {\rm dom}~ g , \bigcup {\rm ran}~
    g , \\
     ( F ` ( g ` \bigcup {\rm dom}~ g ) ) ) ) ) ` ( {\rm rec} ( F , \mc{A} )
    \restriction \mc{B} ) ) )\label{eq:rdgvalg}\tag{rdgvalg}
\end{align}
\end{lemma}


\begin{lemma}\blTitle{rdgval} Value of the recursive definition generator.
(Contributed by NM,
       9-Apr-1995.)  (Revised by Mario Carneiro, 8-Sep-2013.)
\begin{align}
\mm{\vdash} ( \mc{B} \in {\rm On} \rightarrow ( {\rm rec} ( F , \mc{A} ) `
    \mc{B} ) = ( ( g \in {\rm V} \mapsto {\rm if} ( g = \varnothing , \mc{A} ,
    {\rm if} ( {\rm Lim}~ {\rm dom}~ g , \bigcup {\rm ran}~ g , \\
     ( F ` ( g `
    \bigcup {\rm dom}~ g ) ) ) ) ) ` ( {\rm rec} ( F , \mc{A} ) \restriction
    \mc{B} ) ) )\label{eq:rdgval}\tag{rdgval}
\end{align}
\end{lemma}


\begin{lemma}\blTitle{rdg0} The initial value of the recursive definition
generator.  (Contributed
         by NM, 23-Apr-1995.)  (Revised by Mario Carneiro, 14-Nov-2014.)
\begin{align}
\mm{\vdash} \mc{A} \in {\rm V}\quad\Rightarrow\quad \mm{\vdash} ( {\rm rec} ( F
    , \mc{A} ) ` \varnothing ) = \mc{A}\label{eq:rdg0}\tag{rdg0}
\end{align}
\end{lemma}


\begin{lemma}\blTitle{rdgseg} The initial segments of the recursive
definition generator are sets.
       (Contributed by Mario Carneiro, 16-Nov-2014.)
\begin{align}
\mm{\vdash} ( \mc{B} \in {\rm dom}~ {\rm rec} ( F , \mc{A} ) \rightarrow ( {\rm
    rec} ( F , \mc{A} ) \restriction \mc{B} ) \in {\rm V}
    )\label{eq:rdgseg}\tag{rdgseg}
\end{align}
\end{lemma}


\begin{lemma}\blTitle{rdgsucg} The value of the recursive definition
generator at a successor.
       (Contributed by NM, 16-Nov-2014.)
\begin{align}
\mm{\vdash} ( \mc{B} \in {\rm dom}~ {\rm rec} ( F , \mc{A} ) \rightarrow ( {\rm
    rec} ( F , \mc{A} ) ` {\rm suc} \mc{B} ) = ( F ` ( {\rm rec} ( F , \mc{A} )
    ` \mc{B} ) ) )\label{eq:rdgsucg}\tag{rdgsucg}
\end{align}
\end{lemma}


\begin{lemma}\blTitle{rdgsuc} The value of the recursive definition
generator at a successor.
       (Contributed by NM, 23-Apr-1995.)  (Revised by Mario Carneiro,
       14-Nov-2014.)
\begin{align}
\mm{\vdash} ( \mc{B} \in {\rm On} \rightarrow ( {\rm rec} ( F , \mc{A} ) ` {\rm
    suc} \mc{B} ) = ( F ` ( {\rm rec} ( F , \mc{A} ) ` \mc{B} ) )
    )\label{eq:rdgsuc}\tag{rdgsuc}
\end{align}
\end{lemma}


\begin{lemma}\blTitle{rdglimg} The value of the recursive definition
generator at a limit ordinal.
       (Contributed by NM, 16-Nov-2014.)
\begin{align}
\mm{\vdash} ( ( \mc{B} \in {\rm dom}~ {\rm rec} ( F , \mc{A} ) \wedge {\rm Lim}~
    \mc{B} ) \rightarrow ( {\rm rec} ( F , \mc{A} ) ` \mc{B} ) = \bigcup ( {\rm
    rec} ( F , \mc{A} ) `` \mc{B} ) )\label{eq:rdglimg}\tag{rdglimg}
\end{align}
\end{lemma}


\begin{lemma}\blTitle{rdglim} The value of the recursive definition
generator at a limit ordinal.
       (Contributed by NM, 23-Apr-1995.)  (Revised by Mario Carneiro,
       14-Nov-2014.)
\begin{align}
\mm{\vdash} ( ( \mc{B} \in \mc{C} \wedge {\rm Lim}~ \mc{B} ) \rightarrow ( {\rm
    rec} ( F , \mc{A} ) ` \mc{B} ) = \bigcup ( {\rm rec} ( F , \mc{A} ) ``
    \mc{B} ) )\label{eq:rdglim}\tag{rdglim}
\end{align}
\end{lemma}


\begin{lemma}\blTitle{rdg0g} The initial value of the recursive definition
generator.  (Contributed
       by NM, 25-Apr-1995.)
\begin{align}
\mm{\vdash} ( \mc{A} \in \mc{C} \rightarrow ( {\rm rec} ( F , \mc{A} ) `
    \varnothing ) = \mc{A} )\label{eq:rdg0g}\tag{rdg0g}
\end{align}
\end{lemma}


\begin{lemma}\blTitle{rdgsucmptf} The value of the recursive definition
generator at a successor (special
       case where the characteristic function uses the map operation).
       (Contributed by NM, 22-Oct-2003.)  (Revised by Mario Carneiro,
       15-Oct-2016.)
\begin{align}
\mm{\vdash} \underline{\Finv} \msc{x} \mc{A}\quad\&\quad \mm{\vdash}
    \underline{\Finv} \msc{x} \mc{B}\quad\&\quad \mm{\vdash} \underline{\Finv}
    \msc{x} \mc{D}\quad\&\quad \mm{\vdash} F = {\rm rec} ( ( \msc{x} \in {\rm
    V} \mapsto \mc{C} ) , \mc{A} )\quad\&\quad \mm{\vdash} ( \msc{x} = ( F `
    \mc{B} ) \rightarrow \mc{C} = \mc{D} )\quad\Rightarrow\quad \mm{\vdash} ( (
    \mc{B} \in {\rm On} \wedge \mc{D} \in V ) \rightarrow ( F ` {\rm suc}
    \mc{B} ) = \mc{D} )\label{eq:rdgsucmptf}\tag{rdgsucmptf}
\end{align}
\end{lemma}


\begin{lemma}\blTitle{rdgsucmptnf} The value of the recursive definition
generator at a successor (special
       case where the characteristic function is an ordered-pair class
       abstraction and where the mapping class $\mc{D}$ is a proper class). 
This
       is a technical lemma that can be used together with {\tt rdgsucmptf} to
help
       eliminate redundant sethood antecedents.  (Contributed by NM,
       22-Oct-2003.)  (Revised by Mario Carneiro, 15-Oct-2016.)
\begin{align}
\mm{\vdash} \underline{\Finv} \msc{x} \mc{A}\quad\&\quad \mm{\vdash}
    \underline{\Finv} \msc{x} \mc{B}\quad\&\quad \mm{\vdash} \underline{\Finv}
    \msc{x} \mc{D}\quad\&\quad \mm{\vdash} F = {\rm rec} ( ( \msc{x} \in {\rm
    V} \mapsto \mc{C} ) , \mc{A} )\quad\&\quad \mm{\vdash} ( \msc{x} = ( F `
    \mc{B} ) \rightarrow \mc{C} = \mc{D} ) \\
    \quad\Rightarrow\quad \mm{\vdash} (
    \lnot \mc{D} \in {\rm V} \rightarrow ( F ` {\rm suc} \mc{B} ) = \varnothing
    )\label{eq:rdgsucmptnf}\tag{rdgsucmptnf}
\end{align}
\end{lemma}


\begin{lemma}\blTitle{rdgsucmpt2} This version of {\tt rdgsucmpt} avoids
the not-free hypothesis of
       {\tt rdgsucmptf} by using two substitutions instead of one. 
(Contributed by
       Mario Carneiro, 11-Sep-2015.)
\begin{align}
\mm{\vdash} F = {\rm rec} ( ( \msc{x} \in {\rm V} \mapsto \mc{C} ) , \mc{A}
    )\quad\&\quad \mm{\vdash} ( \msc{y} = \msc{x} \rightarrow E = \mc{C}
    )\quad\&\quad \mm{\vdash} ( \msc{y} = ( F ` \mc{B} ) \rightarrow E = \mc{D}
    )\quad\Rightarrow\quad \mm{\vdash} ( ( \mc{B} \in {\rm On} \wedge \mc{D}
    \in V ) \rightarrow ( F ` {\rm suc} \mc{B} ) = \mc{D}
    )\label{eq:rdgsucmpt2}\tag{rdgsucmpt2}
\end{align}
\end{lemma}


\begin{lemma}\blTitle{rdgsucmpt} The value of the recursive definition
generator at a successor (special
       case where the characteristic function uses the map operation).
       (Contributed by Mario Carneiro, 9-Sep-2013.)
\begin{align}
\mm{\vdash} F = {\rm rec} ( ( \msc{x} \in {\rm V} \mapsto \mc{C} ) , \mc{A}
    )\quad\&\quad \mm{\vdash} ( \msc{x} = ( F ` \mc{B} ) \rightarrow \mc{C} =
    \mc{D} )\quad\Rightarrow\quad \mm{\vdash} ( ( \mc{B} \in {\rm On} \wedge
    \mc{D} \in V ) \rightarrow ( F ` {\rm suc} \mc{B} ) = \mc{D}
    )\label{eq:rdgsucmpt}\tag{rdgsucmpt}
\end{align}
\end{lemma}


\begin{lemma}\blTitle{rdglim2} The value of the recursive definition
generator at a limit ordinal, in
       terms of the union of all smaller values.  (Contributed by NM,
       23-Apr-1995.)
\begin{align}
\mm{\vdash} ( ( \mc{B} \in \mc{C} \wedge {\rm Lim}~ \mc{B} ) \rightarrow ( {\rm
    rec} ( F , \mc{A} ) ` \mc{B} ) = \bigcup \{ \msc{y} ~\vert~ \exists \msc{x}
    \in \mc{B} \msc{y} = ( {\rm rec} ( F , \mc{A} ) ` \msc{x} ) \}
    )\label{eq:rdglim2}\tag{rdglim2}
\end{align}
\end{lemma}


\begin{lemma}\blTitle{rdglim2a} The value of the recursive definition
generator at a limit ordinal, in
       terms of indexed union of all smaller values.  (Contributed by NM,
       28-Jun-1998.)
\begin{align}
\mm{\vdash} ( ( \mc{B} \in \mc{C} \wedge {\rm Lim}~ \mc{B} ) \rightarrow ( {\rm
    rec} ( F , \mc{A} ) ` \mc{B} ) = \underline{\bigcup} \msc{x} \in \mc{B} (
    {\rm rec} ( F , \mc{A} ) ` \msc{x} ) )\label{eq:rdglim2a}\tag{rdglim2a}
\end{align}
\end{lemma}


\begin{lemma}\blTitle{rdgprc0} The value of the recursive definition
generator at $\varnothing$ when the base
       value is a proper class.  (Contributed by Scott Fenton, 26-Mar-2014.)
       (Revised by Mario Carneiro, 19-Apr-2014.)
\begin{align}
\mm{\vdash} ( \lnot I \in {\rm V} \rightarrow ( {\rm rec} ( F , I ) `
    \varnothing ) = \varnothing )\label{eq:rdgprc0}\tag{rdgprc0}
\end{align}
\end{lemma}


\begin{lemma}\blTitle{rdgprc} The value of the recursive definition
generator when $I$ is a proper
       class.  (Contributed by Scott Fenton, 26-Mar-2014.)  (Revised by Mario
       Carneiro, 19-Apr-2014.)
\begin{align}
\mm{\vdash} ( \lnot I \in {\rm V} \rightarrow {\rm rec} ( F , I ) = {\rm rec} (
    F , \varnothing ) )\label{eq:rdgprc}\tag{rdgprc}
\end{align}
\end{lemma}


\begin{lemma}\blTitle{rdgsucuni} If an ordinal number has a predecessor,
the value of the recursive
     definition generator at that number in terms of its predecessor.
     (Contributed by ML, 17-Oct-2020.)
\begin{align}
\mm{\vdash} ( ( \mc{B} \in {\rm On} \wedge \mc{B} \ne \varnothing \wedge \lnot
    {\rm Lim}~~ \mc{B} ) \rightarrow ( {\rm rec} ( F , I ) ` \mc{B} ) = ( F ` (
    {\rm rec} ( F , I ) ` \bigcup \mc{B} ) )
    )\label{eq:rdgsucuni}\tag{rdgsucuni}
\end{align}
\end{lemma}


\begin{lemma}\blTitle{rdgeqoa} If a recursive function with an initial
value $\mc{A}$ at step $N$ is
       equal to itself with an initial value $\mc{B}$ at step $M$, then every
       finite number of successor steps will also be equal.  (Contributed by
       ML, 21-Oct-2020.)
\begin{align}
\mm{\vdash} ( ( N \in {\rm On} \wedge M \in {\rm On} \wedge X \in {\rm \omega}
    ) \rightarrow ( ( {\rm rec} ( F , \mc{A} ) ` N ) = ( {\rm rec} ( F , \mc{B}
    ) ` M ) \rightarrow ( {\rm rec} ( F , \mc{A} ) ` ( N +_o X ) ) = ( {\rm
    rec} ( F , \mc{B} ) ` ( M +_o X ) ) ) )\label{eq:rdgeqoa}\tag{rdgeqoa}
\end{align}
\end{lemma}





%


